% !TeX root = ../papers/nc-p3.tex
\section{全ての対称化写像が特異な境界}
\label{sec:singular-transfer-planes}
$Z_{\rm sym}$では、$S$上の四位置間の同一視のグラフを射影行列空間の積で閉じると、
任意の点の上で同時に正証明書$H_{ij}:V_j\to V_i$を得る。
以下ではこのように選んだ一組を固定する。各行列は非零で、
$H_{ik}H_{kj}$と$H_{ij}$は線型従属である。逆向きとの積はスカラーである。

\begin{lemma}[一位置からの三つの写像が二次元の核を共有する場合]\label{lem:three-common-binary-quotient}
一因子階数が全て4の四テンソルで、ある位置$j$から他の三位置への
階数2の正証明書が共通の核$D\subset V_j$を持つとする。
二次関係を隣接対の縮約像で定めると、連続純粋関係が存在する。
\end{lemma}
\begin{proof}
$j$を$B=V_j/D$へ落とす。他の位置$i$は
$A_i=\operatorname{im}H_{ij}$に支持される。実際、各$B\simeq A_i$で
射影した$j$を$i$の空間で読むと、二位置の交換で不変となる。
従って全ての$A_i$を$B$と同一視したテンソルは
三つの互換$(j,i)$で不変であり、$q\in\operatorname{Sym}^4 B$である。
$j$位置の一因子階数は2である。従って$q$の二対二切断の像
$E\subset\operatorname{Sym}^2B$は少なくとも二次元である。
実際、階数1なら対称な切断はスカラーを除き$c\otimes c$で、
中央の二文字の交換から二次形式$c$の二次小行列式は全て零となる。
すると$q$は四乗の倍で、一因子階数に反する。

代数閉体上の射影平面$\mathbb P(\operatorname{Sym}^2 B)$では、
正の次元の線型部分空間は平方の二次曲線に交わる。
よってある$0\ne b\in B$で$b^2\in E$である。
位置$j$以外の三文字を巡回順に見ると、その二つの隣接関係像は、
いずれも対応する同一視のもとで$E$を含む。これは$j$を$D^\perp$で縮約して得る部分像である。
従って$b$に対応する三文字を取れば、その前二文字も後二文字も純粋関係となる。
\end{proof}

\begin{theorem}[対称閉包のAS残余に必要な隣接階数1]\label{thm:singular-transfer-rank-one-required}
$Z_{\rm sym}\setminus S$内の標準AS対象を仮定する。どの同時極限の組にも、
隣接位置間の階数1の正証明書が少なくとも一つある。
\end{theorem}
\begin{proof}
定理\ref{thm:one-positive-AS-in-symmetric-closure}により全ての写像が特異で、
向かい合う位置間では系\ref{cor:opposite-singular-rank-two}から階数2である。
まず全ての向きで階数2と仮定する。各位置$k$で
\[
 A_{ki}=\operatorname{im} H_{ki}=\ker H_{ik}\subset V_k\quad(i\ne k)
\]
と書ける。逆向きとの積が零であることと次元から等号を得る。
$i,j,k$が異なるとき、$H_{ik}H_{kj}$が零なら$A_{ki}=A_{kj}$である。
零でなければ$H_{ij}$に比例し、階数2なのでこの二平面は相補的である。
この場合、合成の像と核から
\[
 A_{ij}=A_{ik},\qquad A_{ji}=A_{jk}
\]
も従う。

ある位置では三平面が全て一致することを示す。
どの位置でも全て一致する場合は明らかである。
そうでなければ位置$k$の三平面を等しいものごとに分け、
単独となる位置$i$を一つ選ぶ。残りの二位置$j,\ell$に対し、
$k$での$i$の平面はそれぞれの平面と相補的なので、
上の合成則は$A_{ij}=A_{ik}=A_{i\ell}$を与える。
位置$i$から出る三写像の核が一致するため、
補題\ref{lem:three-common-binary-quotient}と定理\ref{thm:no-consecutive-pure-relations}が
AS性に矛盾する。従って全階数2は不可能である。

残る写像には階数1か3があり、それは向かい合う位置間にはない。
階数3の場合、同時極限の逆向き写像との積が零であり、逆向きは非零なので階数1である。
従っていずれの場合も同じ組に隣接階数1が存在する。
\end{proof}

三平面の等号関係は各位置で五通りであり、四位置で625通りである。
合成から得た含意を満たすものは29通りで、共通核を持つ位置が二、三、四個のものは
それぞれ12、16、1通りとなる（\path{computations/verify_transfer_plane_patterns.py}）。
この個数は等号関係の組合せに関するものである。

\begin{proposition}[隣接階数1から得る二変数の部分テンソル]\label{prop:forced-binary-endpoint-pencil}
上のAS残余を仮定すると、巡回順またはその逆順で位置を$0,1,2,3$と選び、
位置0を適当な二次元商$B$へ落としたテンソルは、位置2の平面を$B$と同一視して
\[
 x^2\otimes u\otimes v_0+
 (x\otimes y+y\otimes x)\otimes u\otimes v_1+
 y^2\otimes u\otimes v_2
 \quad\in\operatorname{Sym}^2 B_{0,2}\otimes V_1\otimes V_3
\]
と書ける。ここで$x,y$は$B$の基底で、$0\ne u\in V_1$、
$v_0,v_1,v_2\in V_3$は一次独立である。
特に一つの隣接対には像直線が共通の階数1正証明書からなる二次元線型空間が必要である。
\end{proposition}
\begin{proof}
階数1の同時極限を$H_{10}=u\ell$とする。
$\ell$が$\ker H_{20}=\operatorname{im} H_{02}$を消せば
$B=V_0/\ker H_{20}$上の汎関数となる。そうでなければ
$H_{10}H_{02}$は階数1で非零であり、$H_{12}$に比例する。
その共変ベクトルは$\ker H_{02}$を消すから、位置0と2を交換し逆順を取れば
同じ状況に帰着する。$H_{20}$で$B$を位置2の平面と同一視すると、
位置0を$B$へ射影したテンソルは位置0,2の交換で対称である。

$\ell$で縮約した三テンソルは$u\otimes M$で
$M\in B\otimes V_3$である。一因子階数から$M\ne0$であり、
階数1なら連続純粋関係ができるので$\rk M=2$である。
従って$R_1$は$u\otimes B$を含む。$\ell(x)=1,\ell(y)=0$と基底を選べば、
二つの$B$位置における$x^2,xy+yx$の係数は$ku\otimes V_3$に入る。
残る$y^2$の係数を$T\in V_1\otimes V_3$とする。
$T$が$ku\otimes V_3$に入らなければ、$u$を消す汎関数と
位置0の$y$の双対基底で縮約し、$R_2$に$y\otimes v\ne0$を得る。
$u\otimes y\in R_1$に続くので矛盾である。これで表示式を得る。

どの非零の$B^*$による縮約も同じ理由で$M$の階数が2でなければならない。
三係数$v_0,v_1,v_2$が高々二次元しか張らないと、二次元に埋め込んだときの
$2\times2$行列束の行列式は二変数の斉次二次式であり、代数閉体上で射影零点を持つ。
従って三係数は独立である。任意の$\lambda\in B^*$に対し
$u\lambda:V_0\to V_1$は、作用後の二位置をともに$ku$に支持させるので正証明書である。
\end{proof}

この表示で$Z=\langle v_0,v_1,v_2\rangle$とすると、もう一つの向かい合う対について
\[
 u\in\ker H_{31}=\operatorname{im}H_{13},\qquad Z=H_{13}^{-1}(ku)
\]
である。実際、$H_{31}$を位置1に作用させると位置1,3は対称になる。
$H_{31}u\ne0$なら位置3の三次元の支持$Z$が変換後にそのまま残るが、
対称な相手の支持は一次元なので矛盾する。逆方向に$H_{13}$を位置3に作用させると、
位置1の支持が$ku$に入るため$H_{13}Z\subset ku$である。
$\ker H_{13}$は二次元で$\dim Z=3$なので$H_{13}Z=ku$であり、主張を得る。

\begin{theorem}[対称閉包のAS点の完全な認識]\label{thm:symmetric-closure-as-exact}
$Z_{\rm sym}$内の標準AS対象の集合は、対称完全交叉チャート$S$に等しい。
\end{theorem}
\begin{proof}
$S$の全点のAS性は定理\ref{thm:symmetric-frobenius-chart}による。
逆に、$Z_{\rm sym}\setminus S$に標準AS対象があると仮定する。
定理\ref{thm:singular-transfer-rank-one-required}と
命題\ref{prop:forced-binary-endpoint-pencil}の表示を用いる。
そこでの巡回順序の変更は証明の始点の選択であり、逆順の場合は反対代数で議論する。
標準AS性と標準Hilbert級数はこの操作で保たれる。
分類の同値関係に添字移動や反転を加えるものではない。

次の二次元空間と三次元空間を記す。
\[
 \begin{gathered}
 D_0=\ker H_{20},\qquad B_0=V_0/D_0,\qquad
 B_2=\operatorname{im}H_{20},\\
 A_1=\operatorname{im}H_{13},\qquad
 Z=\langle v_0,v_1,v_2\rangle=H_{13}^{-1}(ku).
 \end{gathered}
\]
特に$u\in A_1$であり、命題の表示から
\[
 R_1\supset ku\otimes B_2,\qquad
 \operatorname{im}\bigl(R_0\longrightarrow B_0\otimes V_1\bigr)
       =B_0\otimes ku.                                      \tag{1}
\]
第二の等号は、位置0だけを商にしたテンソルを残る位置2,3で全て縮約して得る。
表示中の三係数$v_0,v_1,v_2$の独立性により、その像は二次元である。

$Z$を消す非零の$\ell\in V_3^*$を取り、位置3で縮約した三テンソルを$c=w(\ell)$とする。
位置3の一因子階数が4なので$c\ne0$である。
位置0を$B_0$へ落とした表示の位置3の支持は$Z$に入るため、
$c\in D_0\otimes V_1\otimes V_2$である。
また$\ker H_{13}\subset Z$より、ある$\eta\in V_1^*$で$\ell=\eta H_{13}$と書ける。
位置3に$H_{13}$を作用させたテンソルは位置1,3の交換で不変であり、
位置3の支持が$A_1$なので位置1の支持も$A_1$に入る。
その位置3を$\eta$で縮約すると
\[
                         c\in D_0\otimes A_1\otimes V_2                 \tag{2}
\]
を得る。

$c$の末端の支持を$E\subset V_2$とし、$e=\dim E$と置く。
このとき$1\le e\le2$である。
実際、位置2で$c$を縮約した像
$J\subset D_0\otimes A_1$の次元は$e$である。
$e\ge3$なら、四次元の$D_0\otimes A_1$の中で
$J$は二次元の$D_0\otimes ku$と非自明に交わる。
従ってある$0\ne d\in D_0$について$d\otimes u\in J\subset R_0$となる。
(1)の$u\otimes B_2\subset R_1$と合わせると連続純粋関係が存在し、
定理\ref{thm:no-consecutive-pure-relations}に反する。

ここで\emph{元の二次代数の商}を作る。
$Q=V_2/E$、$z=\dim Q=4-e\ge2$とし、四位置の一次空間を順に
\[
                        B_0,\quad V_1,\quad Q,\quad k_\ell              \tag{3}
\]
へ写す。最後の写像は$\ell:V_3\to k$である。
$R_2$を$\ell$で縮約した像は$c$の末端支持$E$であり、
$R_3$を$\ell$で縮約した像は$c$の先頭支持で、$D_0$に含まれる。
従って(3)では辺2,3の関係空間全体が零になる。
辺0の関係像は(1)より$B_0\otimes ku$である。
辺1の関係像を$F\subset V_1\otimes Q$とし、
\[
 s=\dim(B_2\cap E),\qquad r=\dim F,\qquad
 t=\dim\bigl(F\cap(ku\otimes Q)\bigr)
\]
と置く。$\dim R_1=6$と(1)により
\[
                             r\le6-s,\qquad t\ge2-s.                    \tag{4}
\]
すなわち、$R_1$の射影の核は$ku\otimes(B_2\cap E)$を含み、
像は$ku\otimes\operatorname{im}(B_2\to Q)$を含む。

三文字ブロックの商
\[
 T=\frac{B_0\otimes V_1\otimes Q}
          {B_0\otimes ku\otimes Q+B_0\otimes F}
\]
の次元$d$は、分母の二部分空間の交叉が
$B_0\otimes(F\cap(ku\otimes Q))$であることから
\[
 d=8z-2z-2r+2t
   \ge 6z-2(6-s)+2(2-s)
   =6z-8\ge4.                                                        \tag{5}
\]
これは射影した最上テンソルの縮約空間の次元ではなく、
各辺の\emph{全二次関係を射影した像}による商の次元である。

辺2,3の関係が零なので、(3)で得た四周期二次商代数の長さ$4m$の経路空間は、
一次元の矢を挿入した$T^{\otimes m}$である。
従って元の代数$C$について$\dim C_{0,4m}\ge d^m\ge4^m$となる。
例えば$m=6$で
\[
               \dim C_{0,24}\ge4096>2925=\binom{27}{3}
\]
となり、標準AS対象のHilbert級数$(1-t)^{-4}$に反する。
これで対称閉包の残余全体を排除した。
\end{proof}
